% Draft of the supplement subsection for study IG.  To use it once experiments/results/ig_summary.json exists:
%   cp docs/paper/ig_section.draft.tex docs/paper/ig_section.tex  (then edit the "Reading" paragraph to what the numbers show).  The include is already in supplement.tex (\IfFileExists).
\subsection{Study IG: The Request Parser on Fresh Operator Texts}\label{sec:moved-ig}
Every bound of the paper is relative to the requested supply, which the parser reads from the operator's words; its prompt was tuned on 40 development utterances and its model chosen on a 60-text held-out set, so no clean test remained. Study IG (registered and stamped before any inference) runs the frozen parser (Qwen3.5-9B, prompt i35) on \IGTexts{} fresh texts written by hand by the assistant that built it: polite, casual, terse, speech-recognition style, English, mixed, indirect needs and negations (accept expected), and two supplies, vague, wrong shelter, unsupported supply, irrelevant and wording traps (refusal expected).
\emph{Claims.} IG1: at most 2 unsafe outcomes (a wrong supply, or an accepted text that should be refused): \IGUnsafe{} (Wilson upper bound \IGUnsafeUpper): \IGOneHolds. IG2: at least 80\% of the legitimate texts accepted with the right supply: \IGAccepted{} (\IGAcceptedRate): \IGTwoHolds. Refusals correct: \IGRefused. Predictions: speech-recognition style \IGAsr{} \IGPredOneHolds; English \IGEnglish{} \IGPredTwoHolds. Other categories (accepted correctly): polite \IGPolite, casual \IGCasual, terse \IGTerse, indirect \IGIndirect, negation \IGNegation.
\emph{Reading.} (Written after the numbers: which texts were unsafe or falsely refused (listed verbatim in \texttt{ig\_summary.json}), what a failed IG1 means for the deployment note's read-back confirmation, and the limits: 136 texts, written by the same assistant, typed imitations of recogniser output, Japanese and English only.)
